\documentclass[11pt]{article}
\usepackage[a4paper,margin=1in]{geometry}
\usepackage{amsmath,amssymb,amsthm}
\usepackage{hyperref}
\usepackage{enumitem}
\usepackage{listings}
\usepackage{xcolor}

\lstdefinelanguage{Lean}{
  morekeywords={namespace,end,def,theorem,structure,abbrev,where,by,import,open,exact,fun,calc,have,show},
  sensitive=true,
  morecomment=[l]{--},
  morecomment=[s]{/-}{-/},
  morestring=[b]"
}

\lstset{
  language=Lean,
  basicstyle=\ttfamily\small,
  keywordstyle=\color{blue!70!black},
  commentstyle=\color{green!40!black},
  stringstyle=\color{red!50!black},
  showstringspaces=false,
  breaklines=true,
  frame=single
}

\title{Lean Source Companion (All Files)}
\author{TheMathBook / Part I / Linear Algebra}
\date{\today}

\newtheorem*{readpoint}{Read Point}

\begin{document}
\maketitle

\section*{Purpose}
This document converts \emph{all Lean files in this project} into a readable study note.
It is not a replacement for Lean; it is a parallel reading layer.

\section*{Target Lean Files}
\begin{itemize}[leftmargin=1.5em]
  \item \texttt{AxlerLean.lean}
  \item \texttt{AxlerLean/AxlerSummary.lean}
  \item \texttt{AxlerLean/Exercises.lean}
  \item \texttt{AxlerLean/Basic.lean}
  \item \texttt{Main.lean}
\end{itemize}

\section{AxlerLean.lean (Library Root)}
\begin{readpoint}
This is the root module. It collects the two core modules:
\texttt{AxlerSummary} and \texttt{Exercises}.
\end{readpoint}

\begin{lstlisting}
import AxlerLean.AxlerSummary
import AxlerLean.Exercises
\end{lstlisting}

\section{AxlerLean/AxlerSummary.lean (Concept + Core Theorems)}

\subsection*{1) Core Data Model}
\begin{itemize}[leftmargin=1.5em]
  \item \texttt{Vec n := Fin n -> Int} : finite-dimensional integer vectors
  \item \texttt{LinearMap n m} : map with additivity + scalar compatibility
\end{itemize}

\begin{lstlisting}
abbrev Vec (n : Nat) := Fin n -> Int

structure LinearMap (n m : Nat) where
  toFun : Vec n -> Vec m
  map_add : ...
  map_smul : ...
\end{lstlisting}

\subsection*{2) Composition Laws}
\begin{itemize}[leftmargin=1.5em]
  \item identity map: \texttt{LinearMap.id}
  \item composition: \texttt{LinearMap.comp}
  \item proved laws:
  \begin{itemize}
    \item \texttt{comp\_assoc}
    \item \texttt{comp\_id\_left}
    \item \texttt{comp\_id\_right}
  \end{itemize}
\end{itemize}

\subsection*{3) Chapter-Aligned Summary (Axler-style flow)}
\paragraph{Ch1 Vector Spaces}
\texttt{zero\_smul\_vec}, \texttt{add\_zero\_vec}, \texttt{neg\_add\_vec}

\paragraph{Ch2 Finite-Dimensional Spaces}
\texttt{vec\_ext}, \texttt{vec\_eq\_of\_all\_coords\_eq}, \texttt{vec\_zero\_unique}

\paragraph{Ch3 Linear Maps}
\texttt{id\_apply}, \texttt{comp\_apply}, \texttt{comp\_preserves\_smul}

\paragraph{Ch4 Polynomials}
\texttt{evalLin}, \texttt{evalLin\_at\_zero}, \texttt{evalLin\_one}

\paragraph{Ch5 Eigenvalues/Eigenvectors}
\texttt{IsEigenvector}, \texttt{eigenvector\_nonzero}, \texttt{eigenvector\_eq\_scale}

\paragraph{Ch6 Inner Product (2D model)}
\texttt{dot2}, \texttt{dot2\_comm}, \texttt{dot2\_zero\_left}

\paragraph{Ch7 Operators}
\texttt{op\_preserves\_add}, \texttt{op\_preserves\_zero}

\paragraph{Ch8 Complex-like model}
\texttt{Pair}, \texttt{conj}, \texttt{conj\_involutive}, \texttt{conj\_re\_fixed}

\paragraph{Ch9 Determinants (2x2 model)}
\texttt{det2}, \texttt{det2\_zero\_of\_equal\_rows}, \texttt{det2\_zero\_matrix}

\section{AxlerLean/Exercises.lean (Problem--Answer Format)}
\begin{readpoint}
This file is the reader-friendly practice layer on top of \texttt{AxlerSummary.lean}.
Problems are tiered:
\texttt{Basic}, \texttt{Standard}, \texttt{Challenge}.
\end{readpoint}

\subsection*{Basic (p1--p6)}
Coordinate identities, identity map action, linear polynomial at zero, zero determinant.

\subsection*{Standard (p7--p12)}
Eigenvector nonzero condition, dot with zero, map zero to zero, conjugation involution,
right identity composition, linear polynomial at one.

\subsection*{Challenge (p13--p18)}
Left identity composition, associativity, zero-vector uniqueness, dot commutativity,
equal-row determinant zero, conjugation fixing real part.

\paragraph{How to read}
\begin{enumerate}[leftmargin=1.5em]
  \item Read problem comment first.
  \item Match theorem name to statement.
  \item Read proof script (\texttt{by ...}) minimally.
  \item Re-run in Lean and modify one line for active recall.
\end{enumerate}

\section{AxlerLean/Basic.lean (Template Artifact)}
\begin{readpoint}
This file is currently a minimal template artifact:
\texttt{def hello := "world"}.
It does not carry chapter math content.
\end{readpoint}

\begin{lstlisting}
def hello := "world"
\end{lstlisting}

\section{Main.lean (Executable Entry)}
\begin{readpoint}
This is the executable entry for \texttt{lake build}. It prints a load message.
\end{readpoint}

\begin{lstlisting}
import AxlerLean

def main : IO Unit :=
  IO.println "AxlerLean: chapter-wise Lean summary loaded."
\end{lstlisting}

\section*{Build}
\subsection*{Lean build}
\begin{verbatim}
cd "TheMathBook/Part I/Linear Algebra"
lake build
\end{verbatim}

\subsection*{LaTeX build}
\begin{verbatim}
cd "TheMathBook/Part I/Linear Algebra/latex"
pdflatex LeanAll.tex
\end{verbatim}

\section*{Recommended Reading Order}
\begin{enumerate}[leftmargin=1.5em]
  \item \texttt{AxlerSummary.lean} (concept map)
  \item \texttt{Exercises.lean} Basic -> Standard -> Challenge
  \item Re-open target theorem in Lean and replay proof
\end{enumerate}

\end{document}
